%% ma: L12-B15-0009
%% lop: 12
%% chuong: 5
%% bai: 15
%% dang: 12.15.7
%% loai: TN4
%% muc: NB
%% dap_an: "C"
%% nguon: "(NBV)-ĐỀ SỐ 14-MỖI NGÀY 1 ĐỀ THI 2025.pdf (D14, MỖI NGÀY 1 ĐỀ - Đề số 14), Phần I Câu 8"
%% kien_thuc: "Bài 15 (phương trình chính tắc của đường thẳng đi qua hai điểm)"
%% trang_thai: chinh-thuc
\begin{ex}
Trong không gian $Oxyz$, cho hai điểm $A(1;-2;3)$ và $B(3;1;1)$. Đường thẳng $AB$ có phương trình là
\choice{$\dfrac{x-1}{4}=\dfrac{y+2}{-1}=\dfrac{z-3}{4}$}{$\dfrac{x-4}{1}=\dfrac{y+1}{-2}=\dfrac{z-4}{3}$}{\True $\dfrac{x-1}{2}=\dfrac{y+2}{3}=\dfrac{z-3}{-2}$}{$\dfrac{x-2}{1}=\dfrac{y-3}{-2}=\dfrac{z+2}{3}$}
\loigiai{$\overrightarrow{AB}=(2;3;-2)$ là vectơ chỉ phương, đường thẳng $AB$ đi qua $A(1;-2;3)$ nên có phương trình $\dfrac{x-1}{2}=\dfrac{y+2}{3}=\dfrac{z-3}{-2}$.}
\end{ex}
